% ================================================================
%  conteudo/cronograma.tex
% ================================================================

\section{Cronograma}

O projeto será desenvolvido ao longo de dois semestres letivos,
compreendendo as etapas de PFC~I e PFC~II, conforme apresentado no
Quadro~\ref{qua:cronograma}.

\quadrocaption{Cronograma de execução do projeto -- PFC I e II}
\label{qua:cronograma}

\begin{table}[H]
\centering
\begin{tabular}{|p{4.8cm}|c|c|c|c|c|c|c|c|c|c|c|c|}
\hline
\textbf{Atividade} &
\textbf{M1} & \textbf{M2} & \textbf{M3} & \textbf{M4} & \textbf{M5} &
\textbf{M6} & \textbf{M7} & \textbf{M8} & \textbf{M9} & \textbf{M10} &
\textbf{M11} & \textbf{M12} \\
\hline
Revisão bibliográfica        & X & X & X & X & X & X & X & X & X & X &   &   \\
\hline
Preparo do inóculo e meios   &   & X & X &   &   &   &   &   &   &   &   &   \\
\hline
Cultivo (15 dias)            &   &   & X & X &   &   &   &   &   &   &   &   \\
\hline
Monitoramento cinético       &   &   & X & X &   &   &   &   &   &   &   &   \\
\hline
Extração de lipídeos         &   &   &   & X & X &   &   &   &   &   &   &   \\
\hline
Análise GC-MS                &   &   &   &   & X & X &   &   &   &   &   &   \\
\hline
Simulação HEFA               &   &   &   &   &   & X & X &   &   &   &   &   \\
\hline
Análise orçamentária         &   &   &   &   &   &   & X & X &   &   &   &   \\
\hline
Análise estatística          &   &   &   &   &   &   &   & X & X &   &   &   \\
\hline
Redação do trabalho          &   &   &   &   & X & X & X & X & X & X &   &   \\
\hline
Apresentação em eventos      &   &   &   &   &   &   &   &   & X & X &   &   \\
\hline
Redação de artigo            &   &   &   &   &   &   &   &   &   & X & X &   \\
\hline
Entrega final (PFC~II)       &   &   &   &   &   &   &   &   &   &   &   & X \\
\hline
\end{tabular}
\end{table}
\fonte{Elaborado pelo autor (2025).}

\noindent Os meses M1--M6 correspondem ao semestre de PFC~I e os meses
M7--M12 ao semestre de PFC~II. Etapas experimentais de maior duração,
como revisão bibliográfica e redação, são conduzidas de forma contínua e
paralela às atividades laboratoriais.

